% !TEX root = ../../mixed-sklyanin-cubics.tex
% Current result: 2026-11-23 / v66.
\subsection{階数一スライスの残余判別曲線と三尖点}
\label{ifq:subsec:caustic}

前小節の $A_2$ 点は孤立した偶然の例ではない。本小節では階数一スライスの
射影中心にある二変数軸を調べ、その軸上の Azumaya 特異点を与える
明示的な判別曲線を求める。これは残余非 smooth locus の
\emph{部分集合}の記述であり、全残余判別を決定するものではない。

引き続き $\Lambda=\lambda E_{11}$、$M=\mu E_{11}$ と置き、
$t=\tau_1$、$u=\tau_2$、$A=\alpha_1^2$、$B=\alpha_2^2$ と書く。
許容条件から $t,u,A,B$ は非零である。
中心 norm は $\mathcal N=-A\det q_x+B\det q_y$ である。
平方変数の射影空間内の直線
\[
 \ell=\{X_2=X_3=Y_2=Y_3=0\}\subset\mathbb P^5
\]
において、式\eqref{ifb:eq:clifford-matrices}から
\begin{equation}\label{ifq:eq:binary}
 4\mathcal N|_\ell=P(X,Y)
 :=AX(\lambda Y-tX)^2-BY(\mu X-uY)^2,\qquad X=X_1,\ Y=Y_1
\end{equation}
となる。さらに $h$ は式\eqref{ifd:eq:ambient-hyperbola}とする。
次の四次多項式を定める：
\begin{equation}\label{ifq:eq:quartic}
 \begin{split}
 Q(\lambda,\mu)={}&4A^2t\lambda^3+AB\lambda^2\mu^2
      +18ABtu\lambda\mu\\
      &-27ABt^2u^2+4B^2u\mu^3.
 \end{split}
\end{equation}

\begin{theorem}[Azumaya 軸に現れる残余非 smooth 曲線]
\label{ifq:thm:axis-discriminant}
二変数三次式の判別式は
\begin{equation}\label{ifq:eq:disc-factor}
 \operatorname{Disc}_X P(X,1)
    =AB(\lambda\mu-tu)^2Q(\lambda,\mu)
\end{equation}
である。特に
\begin{equation}\label{ifq:eq:bad-open}
           Q(\lambda,\mu)=0,\qquad
           \lambda\mu\ne tu,\qquad h(\lambda,\mu)\ne0
\end{equation}
を満たす任意の複素係数では、$\mathcal B$ は六次元 Koszul AS 正則で、
$w$ は正則中心三次であるが、$\mathcal A_{\lambda,\mu}=\mathcal B/(w)$ の
全 coherent 残余圏は非 proper、残余線束圏は非 smooth となる。
より正確には、これらの係数では有限中心射影模型の
Azumaya chart に特異閉点が存在する。

$Q=0$ と $\lambda\mu\ne tu$ の下での軸上の特異点の存在は必要十分である。
ただし、軸外または非 Azumaya 部分の特異点はこの主張では分類しない。
\end{theorem}

\begin{proof}
式\eqref{ifq:eq:binary}を展開すると
\[
 P(X,Y)
 =At^2X^3+(-2At\lambda-B\mu^2)X^2Y
 +(A\lambda^2+2B\mu u)XY^2-Bu^2Y^3
\]
である。三次多項式の標準判別式
$b^2c^2-4ac^3-4b^3d-27a^2d^2+18abcd$
へ四係数を代入して因数分解すれば
式\eqref{ifq:eq:disc-factor}を得る。
$At^2\ne0$、$Bu^2\ne0$ なので重根は $X/Y\in\C^\times$ にある。

軸上では、四つの横変数に関する $\mathcal N$ の一次項は零である。
実際、$X_2=X_3=Y_2=Y_3=0$ で
$\partial_{X_2}\mathcal N=\partial_{X_3}\mathcal N
=\partial_{Y_2}\mathcal N=\partial_{Y_3}\mathcal N=0$ である。
したがって $P$ の重根がちょうど軸上の中心 hypersurface の特異点を与える。
軸上の二つの行列式は
\[
 \det q_x=-\frac{X(\lambda Y-tX)^2}{4},\qquad
 \det q_y=-\frac{Y(\mu X-uY)^2}{4}.
\]
$XY\ne0$ の下でこれらがともに零となるのは
$\lambda\mu=tu$ の場合に限る。中心 norm の零点では
$A\det q_x=B\det q_y$ なので、一方だけが零となることもない。
従って $\lambda\mu\ne tu$ なら、軸上の重根点は
両Clifford形式が非退化な Azumaya 開にある。
式\eqref{ifq:eq:disc-factor}から、同じ開では重根の存在と $Q=0$ は同値である。

$h\ne0$ では命題\ref{ifd:prop:ambient-locus}により周囲代数が
AS 正則となる。$c\ne0$ は同命題で全スライス上確認済みである。
非退化Clifford点では、前小節の
定理\ref{ifd:thm:residual-a2}の証明で行ったように
$Y_1$ と $\omega_x$ を反転し、次数一の中心単元
$\zeta=\omega_x/Y_1$ を用いて次数零の射影chartへ移る。
$w=0$ は $\omega_y=-(\alpha_1/\alpha_2)\omega_x$ を与える。
$Y_1=1$ の座標表示で、その中心の完備局所環は
$P(X,1)$ と横方向の項を合わせた norm hypersurface の
完備局所環である。上の重根ではこの環は特異である。
この chart の代数は中心上次数四の Azumaya 代数であり、
剰余体が $\C$ の完備局所環上で分裂して
$\operatorname{Mat}_4$ になる。

従ってその点の単純加群は、特異な中心局所環の剰余体に
Morita 対応する。特異 noetherian 局所環の剰余体の射影次元が
無限であることは正則局所環の特徴付け
\href{https://stacks.math.columbia.edu/tag/00OC}{Stacks Project, Tag~00OC}
による。極小自由解消の各非零項は極大イデアル係数の微分を持つので、
剰余体の自己 Ext は無限に多くの次数で非零である。
有限中心射影模型上の coherent 一点対象へ延長し、
局所大域 Ext と有限例外 mutation を使えば、
前小節の証明と同じく全 coherent 残余圏が非 proper である。
一方、有限線束圏は proper であるから、smooth であれば
主稿の定理\paperref{main}{crs:thm:coherent}により
coherent 圏への飽和と矛盾する。有限例外 gluing に関する
命題\paperref{main}{crs:prop:residual-smooth}を用いて
残余線束圏の非 smooth 性も従う。
\end{proof}

\begin{remark}[一般点の局所型を検出する横 Hessian]
\label{ifq:rem:transverse}
横変数の次数二までの norm は
\[
 4\mathcal N
 =P(X,Y)-4A C_xX_2X_3+4B C_yY_2Y_3
     +O\bigl((X_2,X_3,Y_2,Y_3)^3\bigr),
\]
ただし
\[
 C_x=X+\frac{t^2}{4}(\lambda Y-tX),\qquad
 C_y=Y+\frac{u^2}{4}(\mu X-uY).
\]
従って $Q=0$ の Azumaya 軸点で $C_xC_y\ne0$ かつ
$P(X,1)$ の重根が正確に二重なら、中心の局所型は
通常二重点 $A_1$ となる。三重根なら横方向の形式的 Morse
消去後に $A_2$ となる。この確認には
$Q$ の重根判別を可換Ku標的の判別式と同一視する必要はない。
\end{remark}

\begin{theorem}[対称因子での三尖点有理四次曲線]
\label{ifq:thm:tricuspidal}
$\tau_1=\tau_2=\tau$、$a_1=a_2\ne0$ とし、
$l=\lambda/\tau$、$m=\mu/\tau$ と置く。
このとき $Q=0$ の被約射影閉包は
\begin{equation}\label{ifq:eq:deltoid}
       l^2m^2+4(l^3+m^3)+18lm-27=0
\end{equation}
という既約有理四次曲線で、その正規化は
\begin{equation}\label{ifq:eq:param}
 \mathbb P^1\longrightarrow\overline{\{Q=0\}},\qquad
 [v:w]\longmapsto
 [-v^4-2vw^3:-2v^3w-w^4:v^2w^2]
\end{equation}
である。$z=v/w\in\C^\times$ の開では
\[
          l=-z^2-\frac2z,\qquad
          m=-2z-\frac1{z^2}.
\]
射影曲線の特異点は正確に三つの通常 cusp
\begin{equation}\label{ifq:eq:three-cusps}
           (l,m)=(-3\xi,-3\xi^2),\qquad \xi^3=1
\end{equation}
である。$z^3\ne\pm1$ の一般点で周囲正則条件が成り立てば、
対応する Azumaya 軸点は $A_1$ となる。
$z^3=1$ の三つの cusp では、
\[
            \tau^3\notin\{1,-8,1/4,-1/2\}
\]
なら周囲AS正則域内の $A_2$ 型 Azumaya 残余特異点を得る。
\end{theorem}

\begin{proof}
$A=B$ なので、式\eqref{ifq:eq:quartic}から
$Q=A^2\tau^4(l^2m^2+4l^3+4m^3+18lm-27)$ を得る。
$f(r)=r(l-r)^2-(mr-1)^2$ と置くと、表示した有理写像上で
\begin{equation}\label{ifq:eq:binary-factor-param}
       f(r)=\frac{(r-z^2)^2(rz^4-1)}{z^4}
\end{equation}
となる。従って正規化写像の像は四次曲線上にある。
式\eqref{ifq:eq:param}は次数四のbasepoint-freeな射であり、
$[1:0:0]$ の逆像は $[1:0]$ の一点、その近傍で
第二座標と第一座標の比は $2(w/v)+O((w/v)^2)$ である。
よって像への一般次数は一となり、像の次数は四である。
表示した四次式はその既約な像を定義する。
同じ結論は、$z^3+lz+2$ と $mz^2-2lz-3$ の
$z$ に関する最後の二つの subresultant が
\[
 (4l^2+lm^2+3m)z+(6l+2m^2),\qquad
 l^2m^2+4l^3+4m^3+18lm-27
\]
となることからも確認できる。

四次式を $\mathcal Q$ と書くと
\[
       l\,\partial_l\mathcal Q-m\,\partial_m\mathcal Q
                           =12(l^3-m^3).
\]
両偏微分が零となる有限点では $l^3=m^3$ である。
$l=m=0$ は曲線上になく、
$l=\xi s,\ m=\xi^2s$、$\xi^3=1$ と置けば
$\partial_l\mathcal Q=2\xi^2s(s+3)^2$ となる。
従って有限特異点は式\eqref{ifq:eq:three-cusps}の三点だけである。
無限遠では $[1:0:0]$ と $[0:1:0]$ しかなく、
$\partial_w\overline{\mathcal Q}$ がそれぞれ $4$ なので滑らかである。
パラメータ表示の一次導関数はともに $z^3-1$ 倍で、
$z^3=1$ では二階・三階導関数が独立である。
ゆえに三点はいずれも通常 cusp で、表示した射は正規化である。

$z^3\ne\pm1$ では式\eqref{ifq:eq:binary-factor-param}の
二重根と単根は相異なり、また
\[
 lm-1=\frac{2(z^3+1)^2}{z^3}\ne0.
\]
従って軸点は Azumaya 開にある。$C_xC_y$ と
$h=(1-\tau^3)^2-\tau^6lm$ は $z$ の非零有理関数なので、
その非消滅開では前注意により通常二重点となる。

$z^3=1$ では軸上の三次式は $(r-z^2)^3$ となり、
$\det q_x=\det q_y=-4\tau^2$、
$C_x=z^2(1-\tau^3)$、$C_y=1-\tau^3$ である。
横 Hessian は非退化であり、軸方向は正確に三次となるから
局所型は $A_2$ である。この三点では
$h=-(4\tau^3-1)(2\tau^3+1)$ なので、
表示した許容条件が周囲正則性を保証する。
\end{proof}

\begin{example}[有理係数の $A_1$ 型残余特異点]
\label{ifq:ex:rational-a1}
$\tau_1=\tau_2=2$、$a_1=a_2=1$、
\[
             \lambda=\frac{15}{2},\qquad \mu=-6
\]
とする。式\eqref{ifq:eq:param}の $z=-1/2$ に対応する。
$Q=0$ かつ $\lambda\mu\ne4$ であり、
\[
 h=769,\quad c=196,\quad
 P(X,1)/196=\frac{(X-16)(4X-1)^2}{4}.
\]
重根 $X=1/4$ において
$\det q_x=\det q_y=-49/16$、$C_x=29/4$、$C_y=-5/2$ である。
$4(\det q_y-\det q_x)$ を五つの局所座標
$(X,X_2,X_3,Y_2,Y_3)$ で微分した Hessian の行列式は
$-10596600\ne0$ となる。
従って射影Azumaya局所代数は
\[
 \operatorname{Mat}_4\!\left(
       \C[[v_1,\ldots,v_5]]/(v_1^2+\cdots+v_5^2)\right)
\]
であり、残余線束圏は非 smooth である。
ここで $P/196$ は重み $A=B=196$ の共通因子を除いた三次式を意味する。
\end{example}

\begin{remark}[全残余判別との区別]
\label{ifq:rem:scope}
$V(Q)\cap D((\lambda\mu-tu)h)$ は、
階数一係数スライスに含まれる残余非 smooth locus の明示部分である。
$V(\lambda\mu-tu)$ は行列式の零となる非 Azumaya 点を含み、
また軸外にも特異点があり得る。これらの有無や全残余判別の等式、
非可換三次と可換Ku標的の方程式の一致、全 smooth 係数への
実比較核の延長は本定理から従わない。
式\eqref{ifq:eq:disc-factor}、横 Hessian、三尖点の
有理表示と有理 $A_1$ 例は
\texttt{computations/verify\_rank\_one\_caustic.py} で
有理係数の多項式恒等式として再検算できる。
\end{remark}
